\label{sec:algorithm}

\subsection{Land Use Character Vector}

The M-track framework (M.0) defines community character as a function
$\mathbf{c}: \mathcal{T} \to \mathbb{R}^k_{\geq 0}$ mapping each census
tract to a non-negative character vector \citep{m0framework}.
For M.2, the character function is the land use fraction vector:
\begin{equation}
  \mathbf{l}(t) = \bigl(
    p_{\mathrm{res}}(t),\;
    p_{\mathrm{com}}(t),\;
    p_{\mathrm{open}}(t),\;
    p_{\mathrm{nat}}(t)
  \bigr)
  \label{eq:landuse}
\end{equation}
where the components are the fractions of NLCD pixels within tract $t$
belonging to each M.2 category.

\paragraph{Component definitions.}
\begin{align}
  p_{\mathrm{res}}(t)  &= \frac{n_{22}(t) + n_{23}(t)}{N(t)} \\
  p_{\mathrm{com}}(t)  &= \frac{n_{24}(t)}{N(t)} \\
  p_{\mathrm{open}}(t) &= \frac{n_{21}(t)}{N(t)} \\
  p_{\mathrm{nat}}(t)  &= \frac{n_{31}(t) + n_{41}(t) + \cdots + n_{95}(t)}{N(t)}
\end{align}
where $n_k(t)$ is the count of class-$k$ pixels within tract $t$,
and $N(t) = n_{21}(t) + n_{22}(t) + n_{23}(t) + n_{24}(t) + n_{31}(t) + \cdots + n_{95}(t)$
is the total non-water pixel count.
Water pixels (class~11) are excluded from the denominator $N(t)$.

The four-component vector sums to 1.0 by construction:
$p_{\mathrm{res}} + p_{\mathrm{com}} + p_{\mathrm{open}} + p_{\mathrm{nat}} = 1$.
This normalization ensures that the cosine similarity between two purely
residential tracts ($\mathbf{l} = (1, 0, 0, 0)$) is 1.0 regardless of
their pixel density, which is the correct behavior: two densely developed
residential neighborhoods are more similar to each other than either is
to a commercial zone, regardless of how many 30-meter pixels each contains.

\paragraph{Water fraction.}
The per-tract water fraction is computed separately:
\begin{equation}
  p_{\mathrm{wat}}(t) = \frac{n_{11}(t)}{n_{11}(t) + N(t)}
  \label{eq:waterfrac}
\end{equation}
and is stored alongside the four-component vector but excluded from the
cosine similarity computation.
It is used only for the hard-cut rule.

\subsection{Cosine Similarity}

The community similarity between adjacent tracts $u$ and $v$ is the cosine
similarity of their land use character vectors:
\begin{equation}
  \mathrm{sim}(\mathbf{l}_u, \mathbf{l}_v)
  = \frac{\mathbf{l}_u \cdot \mathbf{l}_v}{\|\mathbf{l}_u\| \|\mathbf{l}_v\|}
  \label{eq:cosine}
\end{equation}

Table~\ref{tab:landuse-prototypes} shows the character vectors for five
prototypical tract types encountered in American metropolitan areas.

\begin{table}[ht]
\centering
\caption{Prototypical land use character vectors for five tract types.
  All fractions sum to 1.0.}
\label{tab:landuse-prototypes}
\small
\begin{tabular}{lrrrr}
\toprule
Type & $p_{\mathrm{res}}$ & $p_{\mathrm{com}}$ & $p_{\mathrm{open}}$ & $p_{\mathrm{nat}}$ \\
\midrule
Dense residential suburb & 0.75 & 0.03 & 0.18 & 0.04 \\
Strip-mall commercial corridor & 0.05 & 0.85 & 0.08 & 0.02 \\
Industrial park / logistics & 0.02 & 0.88 & 0.08 & 0.02 \\
Exurban open space & 0.12 & 0.01 & 0.35 & 0.52 \\
Rural agricultural & 0.01 & 0.00 & 0.10 & 0.89 \\
\bottomrule
\end{tabular}
\end{table}

\noindent
Note that commercial corridor and industrial park have nearly identical
vectors under the M.2 classification, because NLCD class~24 (high-intensity
developed) does not distinguish commercial from industrial use at 30-meter
resolution.
The distinction between these economic character types is better captured
by M.1 (LODES economic character) than by M.2.
M.2's comparative strength is the residential-to-non-residential transition
and the natural land cover spectrum, not the internal differentiation of
developed zones.

Table~\ref{tab:landuse-similarity} shows the pairwise cosine similarity
between the five prototypical vectors.

\begin{table}[ht]
\centering
\caption{Cosine similarity between prototypical land use tract types.
  Values near 1.0 indicate strong community similarity; values near 0.0
  indicate natural cut seams.}
\label{tab:landuse-similarity}
\small
\begin{tabular}{lrrrrr}
\toprule
 & Dense res. & Commercial & Industrial & Exurban & Rural \\
\midrule
Dense residential & 1.00 & 0.12 & 0.09 & 0.38 & 0.08 \\
Commercial corridor & 0.12 & 1.00 & 0.99 & 0.11 & 0.04 \\
Industrial park & 0.09 & 0.99 & 1.00 & 0.10 & 0.03 \\
Exurban open space & 0.38 & 0.11 & 0.10 & 1.00 & 0.88 \\
Rural agricultural & 0.08 & 0.04 & 0.03 & 0.88 & 1.00 \\
\bottomrule
\end{tabular}
\end{table}

The similarity matrix reveals the geographic community structure:
\begin{itemize}
  \item \textbf{Residential to commercial (0.12).}
    The sharpest community boundary in the M.2 framework is the edge between
    a residential neighborhood and an adjacent commercial corridor.
    This is the natural cut seam M.2 is designed to exploit.
  \item \textbf{Commercial to industrial ($\approx$0.99).}
    NLCD class~24 dominates both; the two tract types are nearly
    indistinguishable from satellite imagery.
    M.1 (economic character) is the appropriate tool for distinguishing these.
  \item \textbf{Rural to exurban (0.88).}
    A soft similarity; both are primarily natural land cover.
    The gradient from agricultural to exurban residential is gentler than
    the urban residential-to-commercial transition.
  \item \textbf{Residential to rural (0.08).}
    A moderate hard boundary between the developed suburban fabric and
    agricultural or forested land: another strong M.2 cut seam.
\end{itemize}

\subsection{The Hard-Cut Rule for Water Boundaries}

The M.2 algorithm imposes a hard override on any adjacency where both
adjacent tracts have majority water classification:
\begin{equation}
  \text{if } \min\!\bigl(p_{\mathrm{wat}}(u),\, p_{\mathrm{wat}}(v)\bigr) > 0.5
  \;\;\Rightarrow\;\; w(u, v) = 0
  \label{eq:hardcut}
\end{equation}

This rule is applied \emph{before} the cosine similarity computation and
takes absolute priority.
Even if two water-adjacent tracts have identical land use vectors, the edge
weight between them is forced to zero, preventing METIS from assigning them
to the same district unless an explicit bridge adjacency overrides the rule.

The threshold of 0.5 is conservative by design.
A tract must be majority water for the hard-cut rule to activate.
Tracts with 30--49\% water pixels (common in coastal areas and along rivers)
do not trigger the hard cut; instead, their water fraction suppresses the
non-water pixel counts, resulting in a naturally lower cosine similarity
with inland tracts.

\paragraph{Bridge adjacency override.}
The \textsc{bisect} adjacency graph can record explicit bridge adjacency edges
in addition to contiguity-derived edges.
For coastal geographies where a bridge or causeway physically connects two
majority-water tracts, the graph constructor can add a bridge edge with a
user-specified weight.
The hard-cut rule applies only to edges derived from polygon contiguity;
bridge edges are not subject to the water cut override.

\subsection{Edge Weight Formula}

The M.2 edge weight follows the M-track blend formula:
\begin{equation}
  w_{\mathrm{lu}}(u, v)
  = \begin{cases}
      0 & \text{if hard-cut rule~\eqref{eq:hardcut} applies} \\[4pt]
      w_{\mathrm{geo}}(u, v) \times
        \bigl(\alpha + (1 - \alpha)\,\mathrm{sim}(\mathbf{l}_u, \mathbf{l}_v)\bigr)
      & \text{otherwise}
    \end{cases}
  \label{eq:weight}
\end{equation}
with default $\alpha = 0.5$.
The blend factor $\alpha$ prevents the weight from collapsing to zero for
dissimilar non-water adjacent tracts, preserving geographic boundary length
as a baseline constraint.

\paragraph{CLI invocation.}
\begin{verbatim}
bisect state --state NC --year 2020 \
    --weights-override land-use \
    --version v_landuse_test
\end{verbatim}
The \texttt{land-use} weight mode reads precomputed NLCD fraction vectors
from the bisect data cache.
Vectors are computed offline using the zonal statistics pipeline described
in Section~\ref{sec:background} and stored in Parquet format alongside
the adjacency graph.

\subsection{Test Invariants}

The following L0 unit test invariants must hold for any correct implementation:

\begin{description}
  \item[\texttt{water\_majority\_gives\_zero\_weight}.]
    For any tract pair $(u, v)$ where $p_{\mathrm{wat}}(u) > 0.5$ and
    $p_{\mathrm{wat}}(v) > 0.5$: $w_{\mathrm{lu}}(u, v) = 0.0$.
  \item[\texttt{identical\_landuse\_gives\_max\_weight}.]
    For any two tracts with identical land use fraction vectors
    $\mathbf{l}_u = \mathbf{l}_v$: $\mathrm{sim}(\mathbf{l}_u, \mathbf{l}_v) = 1.0$
    and $w_{\mathrm{lu}}(u, v) = w_{\mathrm{geo}}(u, v)$ (blend reduces to $\alpha + (1-\alpha) = 1$).
  \item[\texttt{residential\_vs\_commercial\_weight\_low}.]
    For a tract pair where $p_{\mathrm{com}}(u) \geq 0.9$ and
    $p_{\mathrm{res}}(v) \geq 0.9$: $w_{\mathrm{lu}}(u, v) < 0.3 \times w_{\mathrm{geo}}(u, v)$
    (using $\alpha = 0.0$ for this test).
\end{description}
